\documentclass{article}
\usepackage[UTF8]{ctex}
\usepackage{amsmath,amssymb,amsfonts}
\usepackage{enumitem}
\usepackage{geometry}
\geometry{a4paper, margin=1in}
% === 错题本排版优化（不参与题目标识） ===
\setlength{\parindent}{0pt}
\setlength{\parskip}{0.35em plus 0.10em minus 0.08em}
\setlength{\emergencystretch}{3em}
\setlist{itemsep=0.25em, topsep=0.40em, parsep=0pt}
\allowdisplaybreaks[2]
% === 错题本排版优化结束 ===
\begin{document}
\section*{中国科学技术大学2025年数学分析真题}
\section*{105．中国科学技术大学2025年数学分析试卷}

\subsection{题目1}
一．计算题．每题10分，共40分．

\begin{itemize}
  \item[（1）] 计算 \(\lim _{x \rightarrow \infty}\left(\cos \frac{1}{x}\right)^{x^{2}}\) ．
  \item[（2）] 计算 \(\lim _{n \rightarrow \infty} \frac{\sqrt[n]{n!}}{n}\) ．
  \item[（3）] 计算 \(\lim _{n \rightarrow \infty} \int_{0}^{+\infty} x^{2} e^{-x^{2 n}} \mathrm{~d} x\) ．
  \item[（4）] 由方程 \(\sin x+\ln y-x y^{3}=0\) 在 \((0,1)\) 附近确定的隐函数 \(y=f(x)\) ，求 \(f^{\prime}(0)\) 。
\end{itemize}

\subsection{题目2}
二．（15 分）把 \(f(x)=\frac{\pi}{2}-x, x \in[0, \pi]\) 展开成余弦级数，并计算

\[
\sum_{n=1}^{\infty} \frac{1}{(2 n-1)^{2}}, \sum_{n=1}^{\infty} \frac{1}{n^{4}}, \sum_{n=1}^{\infty} \frac{1}{n(n+1)}\left(1+\frac{1}{2}+\cdots+\frac{1}{n}\right) .
\]

\subsection{题目3}
三．（15 分）\(f\) 是定义在 \(D=\left\{(x, y) \mid x^{2}+y^{2} \leq \pi\right\}\) 上的二阶连续可微函数，且

\[
\frac{\partial^{2} f}{\partial x^{2}}+\frac{\partial^{2} f}{\partial y^{2}}=\sin \left(x^{2}+y^{2}\right) .
\]

计算 \(\oint_{\partial D} \frac{\partial f}{\partial n} \mathrm{~d} s\) ，其中 \(n\) 为 \(\partial D\) 的单位外法向量。

\subsection{题目4}
四．（15 分）\(\sum\) 为 \(z=x^{2}+y^{2}, 0 \leq z \leq 1\) ，取下侧，计算

\[
\iint_{\sum}(2 y+z) \mathrm{d} y \mathrm{~d} z+z \mathrm{~d} x \mathrm{~d} y
\]

\subsection{题目5}
五．（15 分）计算 \(\int_{0}^{+\infty} e^{-\alpha x^{2}} \cos (\beta x) \mathrm{d} x\) ，其中 \(\alpha>0, \beta \in \mathbb{R}\) ．

\subsection{题目6}
六．（20 分）定义在 \([0,1]\) 上的连续函数列 \(\left\{f_{n}(x)\right\}_{n \in \mathbb{N}}\) 一致收敛于 \(f(x)\) ．\\
（1）证明：对任意的 \(\varepsilon>0\) ，存在 \(\delta>0\) ，使得当 \(x, y \in[0,1]\) 且 \(|x-y|<\delta\) 时，\(\left|f_{n}(x)-f_{n}(y)\right|<\varepsilon\) 对\(n=1,2, \cdots\) 成立。\\
（2）证明：存在 \(M>0\) ，使得 \(\sup _{x \in[0,1]}\left\{\left|f_{n}(x)\right|\right\} \leq M\) 对所有 \(n=1,2, \cdots\) 成立．

\begin{itemize}
  \item[七．] （20 分）设 \(f(x)=\sum_{n=2}^{\infty} \frac{x^{n}}{n^{2} \ln n}\) ．
  \item[（1）] 证明：\(f(x)\) 在 \([-1,1]\) 上连续．
  \item[（2）] 证明：\(f(x)\) 在 \(x=-1\) 处右可导．
  \item[（3）] 证明： \(\lim _{x \rightarrow 1^{-}} f^{\prime}(x)=+\infty\) ．
  \item[（4）] 证明：\(f(x)\) 在 \(x=1\) 处左不可导．
\end{itemize}

\subsection{题目7}
八．（10 分）证明： \(\mathbb{R}\) 上任意函数的严格极大值点至多可数．
\end{document}
